\documentclass[../main.tex]{subfiles}

\begin{document}

\ifSubfilesClassLoaded{

    \tableofcontents
    
}{}

\chapter{ Measurable Functions }

\section{Elementray Properties of Measruable Functions}

\begin{definition}{}{}
    We endow $\overline{\mathbb{R}}$ with a topology called the order topology in the following manner. For each $a \in \mathbb{R}$ let
\[
L_a = \overline{\mathbb{R}} \cap \{x : x < a\} = [-\infty, a) \quad \text{and} \quad R_a = \overline{\mathbb{R}} \cap \{x : x > a\} = (a, \infty].
\]
The collection $\mathcal{S} = \{L_a : a \in \mathbb{R}\} \cup \{R_a : a \in \mathbb{R}\}$ is taken as a subbasis for this topology. A basis for the topology is given by
\[
\mathcal{S} \cup \{R_a \cap L_b : a, b \in \mathbb{R}, a < b\}.
\]
Observe that the topology on $\mathbb{R}$ induced by the order topology on $\overline{\mathbb{R}}$ is precisely the usual topology on $\mathbb{R}$.
\end{definition}

\begin{definition}{Measurable Functions}{}
   \begin{itemize}
    \item  Suppose $(X, \mathcal{M})$ and $(Y, \mathcal{N})$ are measure spaces. A mapping $f: X \to Y$ is called \dfntxt{measurable with respect to $\mathcal{M}$ and $\mathcal{N}$} if
$f^{-1}(E) \in \mathcal{M} \text{ whenever } E \in \mathcal{N}.$
If there is no danger of confusion, reference to $\mathcal{M}$ and $\mathcal{N}$ will be omitted, and we will simply use the term ``measurable mapping.''
\item If $Y$ is a topological space, a restriction is placed on $\mathcal{N}$. In this case it is always assumed that $\mathcal{N}$ is the $\sigma$-algebra of Borel sets $\mathcal{B}$. Thus, in this situation, a mapping $(X,\mathcal{M}) \xrightarrow{f} (Y,\mathcal{B})$ is measurable if
\[\begin{aligned}
f^{-1}(E) \in \mathcal{M} \text{ whenever } E \in \mathcal{B}.
\end{aligned}\]
\item When $Y$ is taken as $\overline{\mathbb{R}}$ (endowed with the order topology) and $(X,\mathcal{M})$ is a topological space with $\mathcal{M}$ the collection of Borel sets. Then $f$ is called a \dfntxt{Borel measurable function}. 
\item When $X = \mathbb{R}^n$, $\mathcal{M}$ is the class of Lebesgue measurable sets and $Y = \overline{\mathbb{R}}$. Here, the measurability of \(  f  \)  required that $f^{-1}(E)$ be Lebesgue measurable whenever $E \subset \overline{\mathbb{R}}$ is Borel, in which case $f$ is called a \dfntxt{Lebesgue measurable function}. The definitions imply that $E$ is a measurable set if and only if $\chi_E$ is a measurable function.

   \end{itemize}
   
\end{definition}

\begin{theorem}{}{}
    If the mapping $(X, \mathcal{M}) \xrightarrow{f} (Y, \mathcal{B})$ is measurable, where $\mathcal{B}$ is the $\sigma$-algebra of Borel sets.  Define
\[\begin{aligned}
\Sigma = \{E : E \subset Y \text{ and } f^{-1}(E) \in \mathcal{M}\}.
\end{aligned}\]Then \(   \Sigma   \) is a \(   \sigma   \)-algebra.  
\end{theorem}
\begin{proof}
    Easy to verify.
\end{proof}

\begin{corollary}{}{}
    A continuous mapping is a Borel measurable function.
\end{corollary}
\begin{proof}
   If \(  f  \) is continuous, \(   \Sigma   \) contains all open open sets. Since \(   \Sigma   \) is a \(   \sigma   \)-algebra, \(   \Sigma   \) contains all Borel sets.   
\end{proof}

\begin{definition}{}{}
    If \(  f: X\bar{\to}\mathbb{R}   \), we call the sets \[
    \left\{ f> a \right\}:= \left\{ x\in X:f\left( x \right)> a  \right\}
    \]the \dfntxt{superlevel sets of \(  f  \) }. 
\end{definition}

\begin{theorem}{}{}
Let $f: X \to \overline{\mathbb{R}}$, where $(X,\mathcal{M})$ is a measure space. The following conditions are equivalent:
\begin{enumerate}
\item[(i)] $f$ is measurable.
\item[(ii)] $\{f > a\} \in \mathcal{M}$ for each $a \in \mathbb{R}$.
\item[(iii)] $\{f \ge a\} \in \mathcal{M}$ for each $a \in \mathbb{R}$.
\item[(iv)] $\{f < a\} \in \mathcal{M}$ for each $a \in \mathbb{R}$.
\item[(v)] $\{f \le a\} \in \mathcal{M}$ for each $a \in \mathbb{R}$.
\end{enumerate}
\end{theorem}


\begin{theorem}{}{9-3-4}
A function $f: X \to \overline{\mathbb{R}}$ is measurable if and only if
\begin{enumerate}
    \item[(i)] $f^{-1}\{-\infty\} \in \mathcal{M}$ and $f^{-1}\{\infty\} \in \mathcal{M}$ and
    \item[(ii)] $f^{-1}(a,b) \in \mathcal{M}$ for all open intervals $(a,b) \subset \mathbb{R}$.
\end{enumerate}
\end{theorem}


\begin{lemma}{}{}
If $f$ and $g$ are measurable functions, then the following sets are measurable:
\begin{enumerate}
    \item[(i)] $X \cap \{x : f(x) > g(x)\}$,
    \item[(ii)] $X \cap \{x : f(x) \ge g(x)\}$,
    \item[(iii)] $X \cap \{x : f(x) = g(x)\}$.
\end{enumerate}
\end{lemma}

\begin{proof}
    If \(  f\left( x \right)   > g\left( x \right) \), then there is a rational number \(  r  \) such that \(  f\left( x \right)> r> g\left( x \right)    \). Therefore, it follows that \[
    \left\{ f> g \right\}= \bigcup _{r \in \mathbb{Q} }\left( \left\{ f> r \right\} \cap \left\{ g< r \right\}\right), 
    \]which is a measurable set , and easily (i) holds. (ii) is just the complements of (i) with \(  f  \) and \(  g  \) interchanged. (iii) is just the intersection two of two measruable sets of type (ii), and so it too is measurable.     
\end{proof}
\begin{definition}{}{}
    If $f$ and $g$ are real-valued measurable functions, then $f+g$ is undefined at points where it would be of the form $\infty - \infty$. This difficulty is overcome if we define

\[(f+g)(x) := \begin{cases} f(x)+g(x), & x \in X - B, \\ \alpha, & x \in B, \end{cases}\]

where $\alpha \in \overline{\mathbb{R} }$ is chosen arbitrarily and where

\[B := (f^{-1}\{\infty\} \cap g^{-1}\{-\infty\}) \cup (f^{-1}\{-\infty\} \cap g^{-1}\{\infty\}).\]
\end{definition}

\begin{theorem}{}{}
    If \(  f,g  :X\to   \overline{\mathbb{R} } \) are measurable functions, then \(  f+ g  \)  and \(  fg  \) are measurable.  
\end{theorem}
\begin{proof}
We will treat the case that $f$ and $g$ have values in $\mathbb{R}$. The proof is similar in the general case and is left as an exercise (see Exercise 2, Section 5.1).

To prove that the sum is measurable, define $F: X \to \mathbb{R} \times \mathbb{R}$ by
\[
F(x) = (f(x),g(x))
\]
and $G: \mathbb{R} \times \mathbb{R} \to \mathbb{R}$ by
\[
G(x,y) = x + y.
\]
Then $G \circ F(x) = f(x) + g(x)$, so it suffices to show that $G \circ F$ is measurable. Referring to Theorem \ref{thm:9-3-4}, we need only show that $(G \circ F)^{-1}(J) \in \mathcal{M}$ whenever $J \subset \mathbb{R}$ is an open interval. Now $U := G^{-1}(J)$ is an open set in $\mathbb{R}^2$, since $G$ is continuous. Furthermore, $U$ is the union of a countable family, $\mathcal{F}$, of 2-dimensional intervals $I$ of the form $I = I_1 \times I_2$, where $I_1$ and $I_2$ are open intervals in $\mathbb{R}$. Since
\[
F^{-1}(I) = f^{-1}(I_1) \cap g^{-1}(I_2),
\]
we have
\[
F^{-1}(U) = F^{-1}\left(\bigcup_{I \in \mathcal{F}} I\right) = \bigcup_{I \in \mathcal{F}} F^{-1}(I),
\]
which is a measurable set. Thus $G \circ F$ is measurable, since
\[
(G \circ F)^{-1}(J) = F^{-1}(U).
\]
The product is measurable by essentially the same proof.
\end{proof}

\begin{definition}{Cantor Function}{}
    Recall that the Cantor set \(  C  \) can be expressed as \[
    C= \bigcap _{j= 1}^{\infty}C_{j}
    \] Where \(  C_{j}  \) is the union of the \(  2^{j}  \) closed intervals that remain after the \(  j  \)th step of the construction. The set \[
    D_{j}= \left[ 0,1 \right]-C_{j} 
    \]consists of the \(  2^{j-1}  \)    open intervals that are deleted at the \(  j  \)th step. Let there intervals be denoted by \(  I_{j,k}  \), \(  k= 1,2,\cdots ,2^{j-1}  \), and order them in the obvious way from left to right.   Now define a continuous function \(  f_{j}  \) on \(  \left[ 0,1 \right]   \) by \[
    \begin{aligned}
    f_{j}\left( 0 \right)&= 0,\\ 
     f_{j}\left( 1 \right)&= 1,\\ 
      f_{j}\left( x \right)&= \frac{k }{2^{j} },\quad  x \in I_{j,k}  
    \end{aligned}
    \]  and define \(  f_{j}   \) linearly on each interavl of \(  C_{j}  \). The function \(  f_{j}  \) is continuous and nondecreasing, and it satisfies \[
    \left| f_{j}\left( x \right)-f_{j + 1}\left( x \right)   \right|< \frac{1 }{2^{j} },\quad x \in \left[ 0,1 \right]   
    \]   Since \[
    \left| f_{j}-f_{j + m} \right|< \sum _{i= j}^{j + m-1}\frac{1 }{2^{i} }< \frac{1 }{2^{j-1} }   
    \] it follows that the sequence \(  \left\{ f_{j} \right\}  \) converges uniformly to a continuous function \(  f  \), called \dfntxt{Cantor-Lebesgue function}.   
\end{definition}
\begin{note}
    \(  f_{j}  \)在 \(  I_{j,k}  \)上的设定是端点的平均值.  
\end{note}

\begin{corollary}{}{}
The Cantor-Lebesgue function $f$ and its associate $h(x) := f(x) + x$ described above have the following properties:
\begin{enumerate}
    \item[(i)] $f(C) = [0,1]$; that is, $f$ maps a set of measure 0 onto a set of positive measure.
    \item[(ii)] $h$ maps a Lebesgue measurable set onto a nonmeasurable set.
    \item[(iii)] The composition of Lebesgue measurable functions need not be Lebesgue measurable.
    \item[(iv)] Lebesgue measurable function dose not preserve Borel sets. 
\end{enumerate}
\end{corollary}
\begin{proof}
    \begin{enumerate}
        \item \(  f:\left[ 0,1 \right]\to \left[ 0,1 \right]    \) is onto. It is easy to see that \(  f\left( C \right)= \left[ 0,1 \right]    \) , because \(  f\left( C \right)   \) is compact and \(  f\left( \left[ 0,1 \right]-C  \right)   \) is countable.   
        \item 
        \begin{proofsketch}
            把 \(  C  \)看成是集合 \(  \left[ 0,1 \right]   \)测度为零的``骨''. \(  \left[ 0,1 \right]\setminus C   \) 是 \(  \left[ 0,1 \right]   \)的测度为1的``肉''. \(  f  \)的构造是对``骨''的展开: \(  f\left( C \right)= \left[ 0,1 \right]    \); 同时对``肉''的压缩: 令 \(  f  \)在 \(  I_{j,k}  \)上取常值. 通过令\(  h:\left[ 0,1 \right]\to \left[ 0,2 \right]    \).  \(  h\left( x \right)= f\left( x \right)+ x    \), 我们变换映射的尺度, 不再将``肉''压缩, 而是忠实地搬运, 将每个 \(  I_{j,k}  \)映到大小相同的一个区间.       从而间接的看到这两个操作对``骨''的影响区别不大(测度都是1).  
            
        \end{proofsketch}

        \begin{note}
            \(  h\left( C \right)    \)和 \(  f\left( C \right)   \)测度相同在直觉上也是存在依据的.  
            
            
Cantor函数 \(f\) 是一列分段线性函数 \(f_n\) 的一致极限. 在构成Cantor集第 \(n\) 步近似 \(C_n\) 的那些区间上, \(f_n\) 的斜率为 \((3/2)^n\), 这是一个趋向于无穷大的值.

函数 \(h(x) = f(x) + x\) 可以看作是 \(h_n(x) = f_n(x) + x\) 的极限. 在同样的区间上, \(h_n\) 的斜率变为 \((3/2)^n + 1\).

为一个趋于无穷的量加上1, 并不会显著改变它的数量级. 两个斜率的比值为 \(\frac{(3/2)^n + 1}{(3/2)^n} = 1 + (2/3)^n \to 1\).

这便解释了我们的直觉: "+x" 这个操作虽然修复了 \(f\) 对 \(C\) 补集测度的"压缩", 但对于 \(f\) 自身在 \(C\) 上的"测度生成"效应来说, 其*相对*影响是微不足道的. 因此, \(f\) 和 \(h\) 都将测度为零的集合 \(C\) "拉伸"成了一个测度为1的集合.
        \end{note}
        \(  h  \) is trictly increasing. Thus , h is a homeomorphism from \(  \left[ 0,1 \right]   \) onto \(  \left[ 0,2 \right]   \). Then \(  h  \) carries the open set \(  \left[ 0,1 \right]\setminus C   \) to an open set. And we have \[
      \begin{aligned}
        m\left( h\left( \left[ 0,1 \right]\setminus C  \right)  \right)&= m\left( h\left( \bigcup _{j,k}I_{j,k} \right)  \right)   \\ 
         &= \sum _{j,k}m\left( h\left( I_{j,k} \right) \right)\\ 
          &= \sum _{j,k}m\left( I_{j,k} \right) \\ 
          &= 1 
      \end{aligned}
        \] Therefore, \(  h  \) maps the Cantor set  to a set \(  P  \) with measure \(  1  \).  From Exercise \ref{9-7-2} , there exists a non-measurable set \(  N\subseteq P  \). Set \(  A= h^{-1} \left( N \right)   \). Since \(  h^{-1} \left( N \right)\subseteq h^{-1} \left( P \right)= C    \), \(  A  \)  is a subset of measure-zero set \(  C  \).  Thus \(  A  \) is Lebesgue-measurable. But \(  h\left( A \right)= N   \)  is not measurable.
        \item Note that \(  h^{-1}   \) is measurable, since it is continuous. Let \(  F: = h^{-1}   \).  Observe that \(  \chi _{A}  \) is measurable, but \[
        \chi _{A}\circ F= \chi _{A}\circ h^{-1} = \chi _{N}
        \]is not Lebesgue-measurable.
        \item Observe that \(  A  \) is not a Borel set, for if it were, then \(  F^{-1} \left( A \right)   \) would be a Borel set. But \(  F^{-1} \left( A \right)= h\left( A \right)= N    \) and \(  N  \) is not a Borel set. Now \(  g\coloneqq \chi _{A}   \) is a Lebesgue measurable function such that \(  g^{-1} \left( \left\{ 1 \right\} \right)= A   \), where \(  \left\{ 1 \right\}  \) is a Borel sets but \(  A  \) is not.      
    \end{enumerate}
    
\end{proof}

\begin{theorem}{}{}
Suppose $f: X \to \overline{\mathbb{R}}$ is measurable and $g: \overline{\mathbb{R}} \to \overline{\mathbb{R}}$ is Borel measurable. Then $g \circ f$ is measurable. In particular, if $X = \mathbb{R}^n$ and $f$ is Lebesgue measurable, then $g \circ f$ is Lebesgue measurable.
\end{theorem}

\begin{corollary}{}{}
Let $f: X \to \overline{\mathbb{R}}$ be a measurable function.
\begin{enumerate}
    \item[(i)] Let $\varphi(x) = |f(x)|^p$, $0 < p < \infty$, and let $\varphi$ assume arbitrary extended values on the sets $f^{-1}(\infty)$ and $f^{-1}(-\infty)$. Then $\varphi$ is measurable.
    \item[(ii)] Let $\varphi(x) = \frac{1}{f(x)}$, and let $\varphi$ assume arbitrary extended values on the sets $f^{-1}(0), f^{-1}(\infty)$ and $f^{-1}(-\infty)$. Then $\varphi$ is measurable.
\end{enumerate}
In particular, if $X = \mathbb{R}^n$ and $f$ is Lebesgue measurable, then $\varphi$ is Lebesgue measurable in (i) and (ii).

\end{corollary}

\begin{proof}
    \begin{enumerate}
        \item Define \(  g\left( t \right)= \left| t \right|^{p}    \)  for \(  t\in \mathbb{R}   \), and assign arbitrary caluse to \(  g\left( \infty \right)   \) and \(  g\left( -\infty \right)   \). Now apply the prvious theorem.   
        \item Proceed in a similar way by defining \(  g\left( t \right)= \frac{1 }{t }    \)  when \(  t\neq 0  ,\infty,-\infty\) and assigning arebitrary values to \(  g\left( 0 \right)   \), \(  g\left( \infty \right)   \), and \(  g\left( -\infty \right)   \)    .
    \end{enumerate}
    
\end{proof}
\begin{theorem}{}{}
Let $(X, \mathcal{M}, \mu)$ be a complete measure space and let $f, g$ be extended real-valued functions defined on $X$. If $f$ is measurable and $f=g$ almost everywhere, then $g$ is measurable.
\end{theorem}
\begin{remark}
    In a complete measure space, we can consider a function \(  f  \)  that just define on almost everywhere on \(  X  \), by extending \(  f  \) into \(  \overline{f}  \) with arbitrary value taken on where \(  f  \) was not defined.   . 
\end{remark}
\begin{proof}
Let $N = \{x : f(x) = g(x)\}$. Then $\mu(\widetilde{N}) = 0$, and thus $\widetilde{N}$ as well as all subsets of $\widetilde{N}$ are measurable. For $a \in \mathbb{R}$, we have
\[\begin{aligned}
\{g > a\} &= (\{g > a\} \cap N) \cup (\{g > a\} \cap \widetilde{N}) \\
&= (\{f > a\} \cap N) \cup (\{g > a\} \cap \widetilde{N}) \in \mathcal{M}.
\end{aligned}\]
\end{proof}
\begin{theorem}{}{}
Suppose $\varphi$ is an outer measure on a space $X$. Then an extended real-valued function $f$ on $X$ is $\varphi$-measurable if and only if
\[\begin{aligned}
\varphi(A) \geq \varphi(A \cap \{f \leq a\}) + \varphi(A \cap \{f \geq b\})
\end{aligned}\]
whenever $A \subset X$ and $a < b$ are real numbers.
\end{theorem}

\begin{proof}
    The ``only if'' is immediate from \[
     \varphi \left( A \right)\ge  \varphi \left( A\cap \left\{ f\le a \right\} \right)+  \varphi \left( A\cap \left\{ f> a \right\} \right)\ge  \varphi \left( A\cap \left\{ f\le a \right\} \right)+  \varphi \left( A\cap \left\{ f\ge b \right\} \right)     
    \]
    To show the converse, we need to prove that for any \(  r \in \mathbb{R}   \), the set  \[
    E\coloneqq  \left\{ f\le r \right\}
    \] satisfies \[
     \varphi \left( A \right)\ge  \varphi \left( A\cap E \right)+  \varphi \left( A-E \right)   
    \] for each \(  A\subseteq X  \). We define \[
    B_{j}\coloneqq A\cap \left\{ r+   \frac{1 }{j + 1 }\le f\le r+  \frac{1 }{j }   \right\}
    \] We shall show \[
    \infty>  \varphi \left( A \right)\ge  \varphi \left( \bigcup _{j= 1}^{\infty}B_{2j} \right)=   \sum _{j= 1}^{\infty} \varphi \left( B_{2j} \right)  
    \]by induction. We fix \(  k  \) and assume \[
     \varphi \left( \bigcup _{j= 1}^{k}B_{2j} \right)= \sum _{j= 1}^{k} \varphi \left( B_{2j} \right)  
    \] We set \[
    A_{k}\coloneqq \bigcup _{j= 1}^{k}B_{2j}
    \] Then \[
    \left( A_{k}\cup B_{2\left( k+ 1 \right) } \right)\cap \left\{ f\le r+ \frac{1 }{2\left( k+ 1 \right)  }  \right\} = B_{2\left( k+ 1 \right) }
    \] \[
    \left( A_{k}\cup B_{2\left( k+ 1 \right) } \right)\cap \left\{ f\ge r+ \frac{1 }{2k+ 1 }  \right\} = A_{k}
    \]We have \[
     \varphi \left( \bigcup _{j= 1}^{k+ 1}B_{2j} \right)=   \varphi \left( A_{k}\cup B_{2\left( k+ 1 \right) } \right) \ge  \varphi \left( B_{2\left( k+ 1 \right) } \right) +  \varphi \left( A_{k} \right)= \sum _{j= 1}^{k+ 1} \varphi \left( B_{2j} \right)  
    \] Thus we inductively show that \[
    \varphi \left( A \right)\ge    \varphi \left( \bigcup _{j= 1}^{\infty}B_{2j} \right)= \sum _{j= 1}^{\infty} \varphi \left( B_{2j} \right)  
    \] Similarly, one can show that \[
    \varphi \left( A \right)\ge    \varphi \left( \bigcup _{j= 1}^{\infty}B_{2j-1} \right)= \sum _{j= 1}^{\infty} \varphi \left( B_{2j-1} \right)  
    \] Then the series \[
    \sum _{j= 1}^{\infty} \varphi \left( B_{j} \right)\le  2 \varphi \left( A \right)< \infty  
    \]convergents. The tail of the series can be arbitrarily small. For each \(   \varepsilon > 0  \), there exists \(  m \in \mathbb{N}   \) such that \[
     \varepsilon > \sum _{j= m}^{\infty} \varphi \left( B_{j} \right)\ge  \varphi \left( \bigcup _{j= m}^{\infty}B_{j} \right)  
    \]Where \[
    \bigcup _{j = m}^{\infty}B_{j}=  A\cap \left\{ r< f\le r+ \frac{1 }{m }  \right\}
    \] We define a measure \(  \psi   \) \[
    \psi \left( F \right)\coloneqq  \varphi \left( A\cap F \right),\quad \forall F\subseteq X
    \]  Then \[
     \varepsilon > \psi \left( \left\{ r< f\le r+ \frac{1 }{m }  \right\} \right)
    \] Thus \[
    \begin{aligned}
     \varphi \left( A\cap E \right)+  \varphi \left( A-E \right)&= \psi \left( E \right)+ \psi \left( E^{c} \right)\\ 
      &= \psi \left( \left\{ f\le r \right\} \right)+ \psi \left( \left\{ f> r \right\} \right)       \\ 
       &\le  \psi \left( \left\{ f\le r \right\} \right)+ \psi \left( \left\{ f>  r+ \frac{1 }{m } \right\} \right)+ \psi \left( \left\{ r< f\le r+ \frac{1 }{m }  \right\} \right)   \\ 
        &\le \psi \left( \left\{ f\le r \right\} \right) + \psi \left( \left\{ f\ge r+ \frac{1 }{m }  \right\} \right) +  \varepsilon \\ 
         &=  \varphi \left( A\cap \left\{ f\le r \right\} \right)+  \varphi \left( A\cap \left\{ f\ge r+ \frac{1 }{m }  \right\} \right)+  \varepsilon   \\ 
          &\le  \varphi \left( A \right)+  \varepsilon  
    \end{aligned}
    \]Since \(   \varepsilon   \) is arbitrary, we have \[
     \varphi \left( A\cap E \right)+  \varphi \left( A-E \right)\le  \varphi \left( A \right)   
    \] Thus \(  E  \) is measurable. 
\end{proof}

\end{document}